%----------------------------------------------------------------------------------------
% Tailored CV — University Assistant (predoctoral), Soft Matter / Biophysics
% SoMeX Lab, Faculty of Physics, University of Vienna — Job-ID 5570
%----------------------------------------------------------------------------------------

\documentclass[10pt,a4paper,sans]{moderncv}

\usepackage[utf8]{inputenc}
\usepackage[T1]{fontenc}

\moderncvstyle{classic}
\moderncvcolor{blue}

\usepackage[scale=0.78]{geometry}
\usepackage{ragged2e}

%----------------------------------------------------------------------------------------

\firstname{Kundai Farai}
\familyname{Sachikonye}

\address{Auf der Lichtung 19}{82178 Puchheim, Germany}
\mobile{(+49) 1626386344}
\email{kundai.sachikonye@bitspark.com}
\homepage{github.com/fullscreen-triangle}{github.com/fullscreen-triangle}

%----------------------------------------------------------------------------------------

\begin{document}

\makecvtitle

%----------------------------------------------------------------------------------------
%	EDUCATION
%----------------------------------------------------------------------------------------

\section{Education}

\cventry{2019--2021}{PhD candidate, Computational Science}{Technische Universität München / Bayerisches Forschungsinstitut}{}{}{
  First-principles physical modelling, high-dimensional spectral data analysis,
  distributed scientific computing. Departed to pursue independent research.
}
\cventry{2018--2019}{MSc Computational Biology}{Universität Tübingen}{}{}{
  Computational biophysics, statistical physics of biological systems,
  molecular modelling, bioinformatics.
}
\cventry{2014--2017}{MSc Pharmaceutical Biotechnology}{HAW Hamburg}{}{}{
  Physical chemistry, colloidal systems, cell culture, analytical chemistry.
}
\cventry{2011--2014}{BSc Biochemistry and Cell Biology}{Jacobs University Bremen}{}{}{}

%----------------------------------------------------------------------------------------
%	RESEARCH SOFTWARE — SOFT MATTER AND BIOPHYSICS
%----------------------------------------------------------------------------------------

\section{Research Software}

\cvitem{gas-giant}{
  \textbf{First-Principles Fluid Dynamics and Rheology.}
  Derives kinetic gas theory, Navier-Stokes, viscosity, diffusion, and thermal
  conductivity from bounded phase space geometry via Kirchhoff circuit laws ---
  zero empirical parameters.
  Viscosity: $\mu = \tau_c \times g$ (partition lag $\times$ coupling strength);
  yielding transition = phase transition in partition structure (yield stress
  $\sigma_y$ = partition floor of bounded oscillatory system).
  Validation: viscosity 2.9\,\% MAE across 12 pure liquids; Poiseuille flow
  0.9\,\% error; adiabatic expansion 0.1--0.3\,\% error.
  Directly applicable to first-principles prediction of rheological properties
  in yield stress fluids and colloidal suspensions.
  Live: \href{https://gas-giant.vercel.app/simulate}{gas-giant.vercel.app/simulate}.
  \href{https://github.com/fullscreen-triangle/gas-giant}{github.com/fullscreen-triangle/gas-giant}
}

\cvitem{lagrangian}{
  \textbf{First-Principles Optical Physics for DDM.}
  Resonant Stack theorem: Beer-Lambert absorption, Kirchhoff circuit law, and
  thermal emission are formally identical operations ($T=1/e=0.3679$ in all
  three representations). Provides first-principles model of image contrast
  and optical signal formation in DDM: the partition geometry governing bulk
  absorption governs evanescent and scattered light dynamics in the microscope.
  GPU shader telescope: 10/10 benchmarks, $6.9\times10^{-5}$ median error.
  Live: \href{https://space-observatory-xi.vercel.app}{space-observatory-xi.vercel.app}.
  \href{https://github.com/fullscreen-triangle/lagrangian}{github.com/fullscreen-triangle/lagrangian}
}

\cvitem{olduvai}{
  \textbf{Collective Biological Motion Analysis.}
  Axiom $\oint\mathbf{J}\cdot d\mathbf{A}=0$ (zero net current through closed
  surface in biological motion field); derives conservation laws for collective
  dynamics. Rambling/trembling decomposition of motion: supraspinal vs.\ peripheral
  components. Validated on 86 nights wearable sensor data; AUC$=1.00$ for
  Parkinson's/ataxia/healthy classification. Framework extends to bacterial
  collective dynamics: velocity field as current field, topological constraints
  as conservation conditions.
  Live: \href{https://rembling-trambling.vercel.app}{rembling-trambling.vercel.app}.
  \href{https://github.com/fullscreen-triangle/olduvai}{github.com/fullscreen-triangle/olduvai}
}

\cvitem{helicopter}{
  \textbf{Multi-Modal Microscopy Analysis.}
  Sequential constraint exclusion across fluorescence, brightfield, phase-contrast,
  and spectral channels; twelve modalities; sub-nanometre effective resolution.
  Image analysis pipeline applicable to spinning-disk and wide-field DDM data;
  Python-based, GPU-accelerated, browser-deployable.
  Live:
  \href{https://hieronymus-omega.vercel.app/observe/microscopy}{hieronymus-omega.vercel.app}.
  \href{https://github.com/fullscreen-triangle/helicopter}{github.com/fullscreen-triangle/helicopter}
}

\cvitem{borgia}{
  \textbf{Zero-Parameter Molecular Spectroscopy.}
  Electronic transitions, absorption, emission, and vibrational modes from
  bounded phase space geometry. Triple Equivalence: GPU shader = physical
  spectrometer. Applicable to spectral characterisation of colloidal and
  molecular soft matter systems.
  Live: \href{https://honjo-masamune.vercel.app/tools}{honjo-masamune.vercel.app/tools}.
}

%----------------------------------------------------------------------------------------
%	EXPERIENCE
%----------------------------------------------------------------------------------------

\section{Experience}

\cventry{2021--present}{Independent AI \& Computational Physics Developer}{Bitspark GmbH}{}{}{
  First-principles frameworks for fluid dynamics, optical physics, biological
  collective motion, and molecular spectroscopy. All systems validated and
  publicly deployed.
}

\cventry{2019--2021}{Computational Research}{TUM / Bayerisches Forschungsinstitut}{Munich}{}{
  Large-scale spectral data analysis, distributed scientific computing
  (Ray, Dask), federated learning for multi-site data. Python, Java.
}

\cventry{2019}{Process Optimisation}{Fraunhofer Institut IGB}{Stuttgart}{}{
  Real-time sensor data acquisition and ML-based optimisation for industrial
  photobioreactors. Fluid dynamics in biological systems context. Python, C.
}

%----------------------------------------------------------------------------------------
%	TECHNICAL SKILLS
%----------------------------------------------------------------------------------------

\section{Technical Skills}

\cvitem{Physics / Soft Matter}{
  First-principles fluid dynamics (Kirchhoff/Navier-Stokes); rheology theory
  (yield stress, viscosity, viscoelasticity from partition geometry);
  optical physics (Beer-Lambert, scattering, DDM signal formation);
  collective motion conservation laws; thermodynamic phase transitions;
  Banach contraction mappings; spectral analysis (DFT)
}

\cvitem{Programming}{
  Python (primary — NumPy, SciPy, PyTorch, OpenCV, Ray);
  Rust (performance-critical computation);
  JavaScript/TypeScript (WebGPU shaders);
  Bash, SQL; MATLAB (familiar)
}

\cvitem{Image Analysis / Microscopy}{
  Multi-modal microscopy fusion (fluorescence, brightfield, phase-contrast,
  spectral); GPU-accelerated image processing; CNN-based feature extraction;
  instrument-level Python control; large-scale image pipeline construction
}

\cvitem{Infrastructure}{
  Ray, Dask; Docker, FastAPI; WebGPU compute shaders;
  memory-mapped I/O for large datasets; HDF5
}

%----------------------------------------------------------------------------------------
%	LANGUAGES
%----------------------------------------------------------------------------------------

\section{Languages}

\cvitemwithcomment{English}{Native / Fluent (C1+)}{}
\cvitemwithcomment{German}{Conversational}{Living in Germany since 2011}

%----------------------------------------------------------------------------------------

\renewcommand{\listitemsymbol}{-~}

\end{document}
